% Einheit 4 von 8 – Ziehen ohne Zurücklegen und Umlegen (bank/zufallsexperimente-und-pfadregeln/e4.jsonl, zusammenbau v0.1)
%% TODO Zeitmarke Einheit 4: Marken-Zeile nicht lesbar
%% TODO Zweigzeile Teil 1: was hier gelernt wird (Satz in Schülersprache aus dem Einheitstitel) – Einheit 4
\einheitenkopf[e4]{Einheit 4 von 8 · Ziehen ohne Zurücklegen und Umlegen}
\zweigzeile{Abitur GK}
\setcounter{aufgabe}{20}

%% TODO Titel: Ich-kann-Satz fehlt in der Bank (Platzhalter: Kettenname) – Nr. 21
%% TODO Anweisung: ein Satz über der ersten Teilaufgabe fehlt in der Bank – Nr. 21
\begin{aufgabe}{Ohne Zurücklegen und Umlegen}
%% TODO \rechenplatz{n}: Schrittzahl des Grundfalls fehlt in der Bank (nur bei fester Schreibform, 2.3 b)
\begin{teile}
\teil Aus einer Klasse werden nacheinander zwei verschiedene Kinder für zwei Ämter ausgelost. Bleibt P(Mädchen) beim zweiten Losen gleich?\\ \kreuz{p bleibt gleich}\\ \kreuz{p ändert sich}
\teil In einer Klasse mit 25 Kindern tragen 10 eine Brille. Zwei Kinder werden ausgelost. P(beide mit Brille)? P = \leerfeld
\teil Von 40 Befragten sind 16 Vegetarier (V), 24 nicht (N). Die ersten drei werden betrachtet. Vollständiger Baum. A: alle drei Vegetarier. B: mindestens einer Vegetarier. P(A), P(B)?

\baumdrei{V/,N/}{V/,N/}{V/,N/}{V/,N/}{V/,N/}{V/,N/}{V/,N/}
\teil Aus einem Skatblatt (32 Karten, 8 Herz) werden vier Karten gezogen. P(alle vier Herz)? P = \leerfeld
\teil Urne A enthält 3 rote und 1 blaue Kugel, Urne B 2 rote und 2 blaue. Eine Kugel wird zufällig von A nach B gelegt, dann aus B gezogen. P(rot)? P = \leerfeld
\teil In einem Hut liegen 7 Lose, 3 davon Freikarten. Sieben Kinder ziehen nacheinander ohne Zurücklegen. Hat das erste oder das dritte Kind die größere Chance auf eine Freikarte? Begründe. \hfill (Abitur 2021 GK)
\end{teile}
\end{aufgabe}

%% TODO Titel: Ich-kann-Satz fehlt in der Bank (Platzhalter: Kettenname) – Nr. 22
%% TODO Anweisung: ein Satz über der ersten Teilaufgabe fehlt in der Bank – Nr. 22
\begin{aufgabe}{Ohne Zurücklegen und Umlegen – Fehler finden, Begründen, Anwendung, Darstellungswechsel}
\begin{teile}
\teil Urne A: 2 rote, 2 blaue Kugeln; Urne B: 1 rote, 3 blaue. Eine Kugel wird von A nach B gelegt, dann aus B gezogen. Nils: P(rot) $= \frac{1}{4}$. Finde den Fehler.
\teil Begründe: Beim Ziehen ohne Zurücklegen ist die Astwahrscheinlichkeit der zweiten Stufe ein bedingter Anteil.
\teil In einer Kiste mit 24 Akkus sind 2 defekt. Die Kontrolle prüft 3 zufällig gewählte. P(mindestens ein defekter wird entdeckt)? P = \leerfeld
\teil Schreibe zum Baum (ohne Zurücklegen) den Term für P(zwei verschiedene Farben) und berechne ihn.

\baumzwei{R/\frac{3}{7},B/\frac{4}{7}}{R/\frac{2}{6},B/\frac{4}{6}}{R/\frac{3}{6},B/\frac{3}{6}}
\end{teile}
\end{aufgabe}
